\subsection{The odd-prime field budgets for a general target}\label{addo:field}

Fix a primitive target \(\psi\) of order \(o\) and a host of
Proposition~\ref{host:art}. Write \(E\subset\mathbf Q(\mu_N)\) for
its CM field, \(G=\operatorname{Gal}(E/\mathbf Q)\), \(D=|G|\), and
\(c\in G\) for complex conjugation. The exact sequence for \(H_e\) in
that proposition shows that \(D/o\) is a power of two. For an odd
prime \(\ell\), a twist \(e\ge1\), and \(s\in\{1,-1\}\), put
\[
 M_{\ell,e}=S_\ell(E,(\mathbf Q_\ell/\mathbf Z_\ell)(-e)),\qquad
 b_s=(1+sc)/2,\qquad
 h_{\ell,e}^s=\operatorname{length}_{\mathbf Z_\ell} b_s M_{\ell,e}.
\]
Here \(c\) has its natural action on cohomology with the indicated
twist. Use the character weights \(w_{\ell,e,\xi}\) of
Section~\ref{famsel:complement}, with the same definition for either
sign. Write \(F/\mathbf Q_\ell\) for a splitting field for \(G\),
\(O\) for its ring of integers, and \(a_O=e(O/\mathbf Z_\ell)\).

\begin{proposition}[Comparison on each sign]\label{addo:compare}
For either sign and every odd prime,
\[
 (\log\ell)h_{\ell,e}^s
 \ \le\ \sum_{\xi(c)=s} w_{\ell,e,\xi}
       + C_\psi v_\ell(o)(\log\ell)D(1+\log(2N)).
\]
If \(\ell\nmid o\), restriction gives equality with the character sum.
The total error on any finite set of odd primes and all \(1\le e<t\)
is \(O_\psi(tD(1+\log(2N)))\).
\end{proposition}

\begin{proof}
Put \(n=e+1\), \(U=\operatorname{Spec}\mathcal O_E[1/(N\ell)]\), and
\(H=H^2_{\rm et}(U,\mathbf Z_\ell(n))\). At every finite
\(w\nmid\ell\), an unramified local class with the field coefficients
is zero. Indeed its group is the quotient of the divisible coefficient
module by the image of \((\mathrm Nw)^{-e}-1\), a nonzero scalar.
At \(w\mid\ell\) local vanishing is imposed. Thus the strict field
group is the kernel of localization for the restricted Galois group
with the primes over \(N\ell\) allowed. Poitou--Tate duality and the
degree-two comparison in the proof of Theorem~\ref{newcoh:lemmaK}
give a \(G\)-equivariant injection \(M_{\ell,e}^{\vee}\hookrightarrow H\).
The group \(H\) is finite. Over \(\mathcal O_E[1/\ell]\) this follows
from the even Quillen--Lichtenbaum comparison and finiteness of
\(K_{2e}(\mathcal O_E)\) as in Proposition~\ref{newcoh:fullsupport}.
Further localization adds subquotients of
\(H^1(k_w,\mathbf Z_\ell(e))\), finite since \((\mathrm Nw)^e\ne1\).
The character groups are finite by Proposition~\ref{newcoh:fullsupport}.

For \(\ell\mid o\) we need a bound for the number of generators,
uniform in the twist. The coefficient exact sequence injects
\(H/\ell H\) into \(H^2_{\rm et}(U,\mu_\ell^{\otimes n})\). Restrict
to \(L=E(\mu_\ell)\), writing \(d_L=[L:\mathbf Q]\) and
\(U_L=\operatorname{Spec}\mathcal O_L[1/(N\ell)]\). This restriction
is injective since the cover is finite \'etale of degree dividing
\(\ell-1\) and trace composes with restriction to multiplication by its
degree. A primitive \(\ell\)-th root identifies the coefficient sheaf
on \(U_L\) with \(\mu_\ell\). The Kummer sequence gives
\[
 \dim_{\mathbf F_\ell} H/\ell H
 \le \dim_{\mathbf F_\ell}\operatorname{Pic}(U_L)/\ell
       +\dim_{\mathbf F_\ell}\operatorname{Br}(U_L)[\ell].
\]
The first term is at most \(\log h_L/\log\ell\), where \(h_L\) is
the class number. Minkowski gives each ideal class an integral ideal
of norm at most \(2^{d_L}|\Delta_L|^{1/2}\). The number of ideals
of norm at most \(y\) is at most \(y^2\zeta_L(2)\le y^2\zeta(2)^{d_L}\).
Hence \(\log h_L\ll d_L+\log|\Delta_L|\). The cyclotomic discriminant
bound and the discriminant formula in towers give
\(\log|\Delta_L|\le d_L\log(\ell N)\), since \(L\subset\mathbf Q(\mu_{\ell N})\).
By local invariants the Brauer term is at most the number of places
of \(L\) over \(N\ell\), which is at most \(d_L(1+\omega(N))\).
Here \(L\) is totally imaginary, and \(\omega(N)\) counts the distinct
prime divisors of \(N\). Since \(d_L\le(\ell-1)D\) and \(\ell\mid o\),
we obtain
\[
 r_\ell:=\dim_{\mathbf F_\ell} H/\ell H
       \ll_\psi D(1+\log(2N)).
\]
Taking \(\ell\)-torsion in the injection of \(M_{\ell,e}^{\vee}\)
and using the invariant factors of the finite group \(H\) shows that
\(M_{\ell,e}\), and hence each direct summand \(b_s M_{\ell,e}\),
has at most \(r_\ell\) generators.

Let \(M_O=M_{\ell,e}\otimes_{\mathbf Z_\ell}O\). For a host character
\(\xi\) put \(T_\xi=\sum_{g\in G}\xi(g)^{-1}g\in O[G]\).
A class in \(M_O\) can be viewed over \(E\) with coefficients
\(A_\xi=(F/O)(\xi^{-1}\chi_\ell^{-e})\), since \(\xi\) is trivial on
\(G_E\). Finite free extension of scalars commutes with the cochains
and the local kernels. The class vanishes locally at every finite
place of \(E\), so its corestriction vanishes locally at every finite
place of \(\mathbf Q\) by the local restriction--corestriction formula.
In particular this corestriction lies in \(S_{\ell N}(\mathbf Q,A_\xi)\).
After untwisting over \(E\), its restriction is \(T_\xi\) applied to
the original class. Thus \(T_\xi M_O\) lies in the restriction image
of that strict character group. Character orthogonality gives the
integral identity \(\sum_{\xi(c)=s}T_\xi=D b_s\). Consequently the
length of \(D b_s M_O\) is bounded by the sum of the character lengths
of sign \(s\). At \(\ell\mid o\) the quotient \(b_s M_O/D b_s M_O\)
has length at most \(a_Ov_\ell(D)r_\ell\). Dividing by \(a_O\) and
multiplying by \(\log\ell\) proves the stated inequality. At
\(\ell\nmid o\) the exact restriction descent of
Proposition~\ref{zetafield:descent} applies verbatim, since its
averaging and local restriction arguments require only \(\ell\nmid D\).
Finally \(\sum_{\ell\mid o}v_\ell(o)\log\ell\le\log o\).
\end{proof}

\begin{theorem}[Uniform opposite-parity field sum]\label{addo:opposite}
For every finite set \(P\) of odd primes and every integer \(t\ge2\),
\[
 \sum_{\ell\in P}(\log\ell)\sum_{1\le e<t} h_{\ell,e}^{(-1)^{e+1}}
     \ \ll\ D t^2\bigl(D+\log(2Nt)\bigr)
\]
with an absolute constant.
\end{theorem}

\begin{proof}
Retain \(n=e+1\), \(U\), and \(H\) from the preceding proof. Set
\(b=b_{(-1)^n}\), \(R=b\mathbf Z_\ell[G]\), and let \(\Sigma\) contain
the infinite place and precisely the finite primes dividing \(N\).
By the refined Coates--Sinnott theorem
\cite[Theorem~6.11 and Lemma~6.12]{greitherPopescuCS},
the \(b\)-part of the zeroth Fitting ideal of \(H\) is the
annihilator of \(H^1_{\rm et}(U,\mathbf Z_\ell(n))_{\rm tors}\)
times the equivariant value \(\Theta_\Sigma(1-n)\). Here \(E/\mathbf Q\)
is abelian CM, \(\ell\) is odd, and \(\Sigma\) contains all ramified
primes. The required cyclotomic \(\mu\)-invariant vanishes by
Ferrero--Washington~\cite{ferreroWashington} for the abelian field
\(E(\mu_\ell)\). The cited lemma permits \(\ell\notin\Sigma\).
As in Theorem~\ref{zetafield:minus},
\[
 W_n=(3^{Dn}-1)(5^{Dn}-1),\qquad C_n=nN(n+1)!.
\]
The first integer annihilates the torsion in that \(H^1\). By the
coefficient exact sequence the torsion classes come from the
invariants of \((\mathbf Q_\ell/\mathbf Z_\ell)(n)\),
on which Frobenius over \(p\in\{3,5\}\), \(p\ne\ell\), acts by
\(p^{f_wn}\) with \(f_w\mid D\). In particular
\[
 x=C_nW_n b\Theta_\Sigma(1-n)\ \in\ \operatorname{Fitt}_R(bH).
\]
For every character of sign \((-1)^n\) its value is
\(x_\xi=C_nW_n L_\Sigma(1-n,\xi^{-1})\). These values are nonzero
by the functional equation and the nonvanishing of the removed Euler
factors. The generalized Bernoulli formula at modulus \(N\) gives
\(C_n L_\Sigma(1-n,\xi^{-1})\) integral and
\(\lvert L_\Sigma(1-n,\xi^{-1})\rvert\le 4n!N^n\) at every complex
embedding, exactly as in the proof of Theorem~\ref{zetafield:minus}.
Consequently their Galois-stable product satisfies
\[
 P_n=\prod_{\xi(c)=(-1)^n} x_\xi\in\mathbf Z\setminus\{0\},
 \qquad \log|P_n|\ll Dn\bigl(D+\log(2Nn)\bigr).
\]

We show directly that \(\operatorname{length}_{\mathbf Z_\ell} bH
\le v_\ell(P_n)\), including at \(\ell\mid D\). Take an \(R\)-presentation
\(R^m\xrightarrow{A}R^k\to bH\to0\). We may suppose \(k>0\).
Since \(bH\) is finite, each character matrix \(A_\xi\) of this sign
has full row rank in \(F\), so \(m\ge k\). With \(v_O\) the normalized
valuation and \(J\) ranging over the \(k\)-element column sets, put
\(u_\xi=\min_J v_O(\det A_{\xi,J})\). The Fitting ideal is generated
by the maximal minors, so \(u_\xi\le v_O(x_\xi)\).

For an \(m\)-by-\(k\) matrix of indeterminates \(Y\), Cauchy--Binet
shows that \(\pi_O^{-u_\xi}\det(A_\xi Y)\) is integral and has nonzero
reduction. Indeed the monomials in \(\det Y_J\) use exactly the rows
in \(J\), and at least one such minor has a unit coefficient after
normalization. The product of these finitely many reduced polynomials
is nonzero over an algebraic closure of the residue field. Choose a
point at which it is nonzero and lift its coordinates to a finite
unramified extension \(O'/O\). The resulting single scalar matrix
\(Y\) satisfies \(v_{O'}(\det(A_\xi Y))=u_\xi\) for every \(\xi\).
Let \(R'=R\otimes_{\mathbf Z_\ell}O'\). The cokernel of
\(AY:(R')^k\to(R')^k\) surjects onto \(bH\otimes_{\mathbf Z_\ell}O'\).
The module \(R'\) is finite free over \(O'\), and over its fraction
field the sign algebra splits as one copy for every \(\xi(c)=(-1)^n\).
The determinant of \(AY\) as an \(O'\)-linear endomorphism is therefore
\(\prod_\xi\det(A_\xi Y)\). Taking lengths gives
\[
 a_O\operatorname{length}_{\mathbf Z_\ell} bH
 \le \sum_\xi u_\xi\le \sum_\xi v_O(x_\xi)=a_Ov_\ell(P_n).
\]
Duality preserves the sign of the involution \(c\), so the injection
\(M_{\ell,e}^{\vee}\hookrightarrow H\) gives the same upper bound for
\(h_{\ell,e}^{(-1)^n}\). Summing the prime valuations costs at most
\(\log|P_n|\), and summing over \(2\le n\le t\) proves the estimate.
\end{proof}

\begin{corollary}[The residual Selmer sum]\label{addo:remainder}
Fix the target and \(\varepsilon\) as in Theorem~\ref{adds:dyadic-window}.
On its jointly selected late interval, take \(N,E,t=t(v)\) at the
same page. For any finite set \(S\) of primes, start with the triples
\((\ell,e,\xi)\) with \(\ell\in S\) odd, \(1\le e<t\), and
\(\xi(c)=(-1)^e\). Delete Cases IV and Z, the conductor range of
Corollary~\ref{famsel:subwindow}, and the three order ranges of
Proposition~\ref{adds:order}. Denote the rest by \(\mathcal R(S,v)\).
Then
\[
 \sum_{\ell\in S}(\log\ell)\sum_{1\le e<t} h_\ell(E,e)
 \ \le\ \sum_{(\ell,e,\xi)\in\mathcal R(S,v)}w_{\ell,e,\xi}
        +O_{\psi,\varepsilon}\!\left(\exp(v^{4\varepsilon/5})\right).
\]
One may also delete entries proved zero by
Propositions~\ref{famsel:class}, \ref{famsel:primary} and~\ref{famsel:splitray}.
A bound \(o(v^\varepsilon)\) for the logarithm
of two plus the residual sum, uniform in \(S,v\) on these windows, suffices
for the entire field-length budget of Formal input~\ref{red:design}.
\end{corollary}
\begin{proof}
Prescribe the fixed fields of Corollary~\ref{adds:fixed-char} in
Theorem~\ref{adds:dyadic-window} when choosing the common interval.
The latter theorem bounds the term at two. Apply Theorem~\ref{addo:opposite}
to the opposite sign and Proposition~\ref{addo:compare} to the matching
sign. Their field bound and total comparison error are polynomial in
\(D,t,\log(2Nt)\), with \(D\le N\) and \(\log(Nt)\le 2v^{\varepsilon/4}\).
The fixed fields bound Case IV. Case Z vanishes by
Proposition~\ref{newcoh:zero} since \(3\in B_\psi\).
In the small-conductor range the proof of Corollary~\ref{famsel:subwindow}
gives \(v^{\varepsilon/102}\exp(v^{\varepsilon/2})\).
Proposition~\ref{adds:order} bounds the union of the three order
ranges. The weights are nonnegative, so overlaps in these upper bounds
cost no more than their sum. The asserted logarithmic condition now
makes the full sum at most \(\exp(v^\varepsilon/2)\) eventually.
The rescaling in Section~\ref{sup:section} keeps these same hosts.
\end{proof}
